\lysh{14744}{4}

\begin{alist}
\item $x - x^2 \le \dfrac{1}{4} \Leftrightarrow x^2 - x + \dfrac{1}{4} \ge 0 \Leftrightarrow x^2 - 2 \cdot x \cdot \dfrac{1}{2} + \left(\dfrac{1}{2}\right)^2 \ge 0 \Leftrightarrow \left(x - \dfrac{1}{2}\right)^2 \ge 0$, ισχύει.
Το ίσον ισχύει αν και μόνον αν $x - \dfrac{1}{2} = 0 \Leftrightarrow x = \dfrac{1}{2}$.

\item 
\begin{rlist}
\item Αρχικά παρατηρούμε ότι το σημείο $\Delta$ έχει τετμημένη μηδέν και τεταγμένη $f(0) = 1$,
ενώ αν είναι $k$ η τετμημένη του σημείου $E$, η τεταγμένη του θα είναι $f(k) = 0 = 1 - k$, άρα $k = 1$. Ώστε $\Delta(0, 1)$ και $E(1, 0)$.
Έστω τώρα ότι $\alpha$ είναι η τετμημένη του σημείου $A$, με $0 \le \alpha \le 1$, οπότε το σημείο $A$
έχει τεταγμένη $f(\alpha) = 1 - \alpha$.
Έτσι έχουμε $\Gamma(\alpha, 0), A(\alpha, 1 - \alpha), B(0, 1 - \alpha)$.
Επομένως το εμβαδόν του ορθογωνίου $ABO\Gamma$ είναι $E = \alpha \cdot (1 - \alpha) = \alpha - \alpha^2$.

\item Σύμφωνα με το $\alpha)$ ερώτημα είναι $E \le \dfrac{1}{4}$, δηλαδή για οποιαδήποτε τιμή του $\alpha$, το
εμβαδόν του ορθογωνίου γίνεται το πολύ $\dfrac{1}{4}$. Το ίσον ισχύει αν και μόνον αν $\alpha = \dfrac{1}{2}$,
οπότε σε αυτή την περίπτωση είναι $A\left(\dfrac{1}{2}, \dfrac{1}{2}\right)$, δηλαδή το $ABO\Gamma$ θα γίνει τετράγωνο.
\end{rlist}
\end{alist}

